% GENERATED FILE. Edit canonical lessons, then run npm run generate:books.
% canonical-source-hash: fnv1a64:6082b696c8f82c74
% canonical-lessons: FR-C231-retour, FR-C231-arrivee, FR-C231-depart, FR-C231-destination, FR-C231-paysage

\chapter{Return, Arrival, Departure, Destination, Landscape}
\label{ch:fr-return-arrival-departure-destination-landscape}
\hlchaptermodality{231}

\begin{chapteropening}
\textbf{By the end of this chapter:} \emph{I can say the return, arrival, departure, a destination and a landscape.}

Complete the last lesson of chapter 231: I can say the return, arrival, departure, a destination and a landscape.
\end{chapteropening}

\section[le retour]{le retour — the return}

\label{lesson:FR-C231-retour}



\begin{quote}
Before the new one: say the French for a hammer, then the French for a ladder.
\end{quote}

\subsection*{You'll want to know: le retour}
\textbf{le retour} — \textquotedblleft{}the return\textquotedblright{}.

\begin{cousinweb}
\textbf{retour} — formed from the verb for \textquotedblleft{}to return\textquotedblright{}, \emph{re-} with \emph{tourner}, \textquotedblleft{}to turn\textquotedblright{}. \emph{Tourner} $\leftarrow$ Latin \emph{tornāre}, \textquotedblleft{}to turn\textquotedblright{} (EtymDB 2.1).
\end{cousinweb}

\begin{grammarlens}[title={What you've built}]
One more word for this chapter.
\end{grammarlens}

\subsection*{Your turn}
\begin{itemize}
\raggedright
  \item \emph{Say it:} \emph{le retour}
  \item \emph{Say it:} \emph{le retour}, once more
  \item \emph{Say it:} say \emph{le retour}
  \item \emph{Recall:} say the French for a hammer, then the French for a ladder, then say \emph{le retour} again
\end{itemize}

\begin{tcolorbox}[breakable,colback=teal!4,colframe=teal!35!black,title={Before you move on}]
What does le retour mean? (\textquotedblleft{}The return\textquotedblright{}.) Say it once more.
\end{tcolorbox}
\section[l'arrivée]{l'arrivée — arrival}

\label{lesson:FR-C231-arrivee}



\begin{quote}
Before the new one: say the French for a ladder, then the French for the return.
\end{quote}

\subsection*{You'll want to know: l'arrivée}
\textbf{l'arrivée} — \textquotedblleft{}arrival\textquotedblright{}.

\begin{cousinweb}
\textbf{arrivée} is the past participle of the verb for \textquotedblleft{}to arrive\textquotedblright{}, $\leftarrow$ Vulgar Latin \emph{*arripare}, \textquotedblleft{}to reach the shore\textquotedblright{}, from \emph{ripa}, \textquotedblleft{}a bank\textquotedblright{}, the Latin of \emph{rivière}.
\end{cousinweb}

\begin{grammarlens}[title={What you've built}]
One more word for this chapter.
\end{grammarlens}

\subsection*{Your turn}
\begin{itemize}
\raggedright
  \item \emph{Say it:} \emph{l'arrivée}
  \item \emph{Say it:} \emph{l'arrivée}, once more
  \item \emph{Say it:} say \emph{l'arrivée}
  \item \emph{Recall:} say the French for a ladder, then the French for the return, then say \emph{l'arrivée} again
\end{itemize}

\begin{tcolorbox}[breakable,colback=teal!4,colframe=teal!35!black,title={Before you move on}]
What does l'arrivée mean? (\textquotedblleft{}Arrival\textquotedblright{}.) Say it once more.
\end{tcolorbox}
\section[le départ]{le départ — departure}

\label{lesson:FR-C231-depart}



\begin{quote}
Before the new one: say the French for the return, then the French for arrival.
\end{quote}

\subsection*{You'll want to know: le départ}
\textbf{le départ} — \textquotedblleft{}departure\textquotedblright{}.

\begin{cousinweb}
\textbf{départ} — formed from Old French \emph{departir}, \textquotedblleft{}to divide, to go away\textquotedblright{}, from \emph{de-} and Old French \emph{partir}, \textquotedblleft{}to share out\textquotedblright{}, $\leftarrow$ Latin \emph{partire}.
\end{cousinweb}

\begin{grammarlens}[title={What you've built}]
One more word for this chapter.
\end{grammarlens}

\subsection*{Your turn}
\begin{itemize}
\raggedright
  \item \emph{Say it:} \emph{le départ}
  \item \emph{Say it:} \emph{le départ}, once more
  \item \emph{Say it:} say \emph{le départ}
  \item \emph{Recall:} say the French for the return, then the French for arrival, then say \emph{le départ} again
\end{itemize}

\begin{tcolorbox}[breakable,colback=teal!4,colframe=teal!35!black,title={Before you move on}]
What does le départ mean? (\textquotedblleft{}Departure\textquotedblright{}.) Say it once more.
\end{tcolorbox}
\section[la destination]{la destination — a destination}

\label{lesson:FR-C231-destination}



\begin{quote}
Before the new one: say the French for arrival, then the French for departure.
\end{quote}

\subsection*{You'll want to know: la destination}
\textbf{la destination} — \textquotedblleft{}a destination\textquotedblright{}.

\begin{cousinweb}
\textbf{destination} $\leftarrow$ Latin \emph{destinatio}, \textquotedblleft{}a fixing, a purpose\textquotedblright{}, from \emph{destinare}, \textquotedblleft{}to make fast\textquotedblright{} (Lewis \& Short) — a learned borrowing.
\end{cousinweb}

\begin{grammarlens}[title={What you've built}]
One more word for this chapter.
\end{grammarlens}

\subsection*{Your turn}
\begin{itemize}
\raggedright
  \item \emph{Say it:} \emph{la destination}
  \item \emph{Say it:} \emph{la destination}, once more
  \item \emph{Say it:} say \emph{la destination}
  \item \emph{Recall:} say the French for arrival, then the French for departure, then say \emph{la destination} again
\end{itemize}

\begin{tcolorbox}[breakable,colback=teal!4,colframe=teal!35!black,title={Before you move on}]
What does la destination mean? (\textquotedblleft{}A destination\textquotedblright{}.) Say it once more.
\end{tcolorbox}
\section[le paysage]{le paysage — a landscape}

\label{lesson:FR-C231-paysage}



\begin{quote}
Before the new one: say the French for departure, then the French for a destination.
\end{quote}

\subsection*{You'll want to know: le paysage}
\textbf{le paysage} — \textquotedblleft{}a landscape\textquotedblright{}.

\begin{cousinweb}
\textbf{paysage} — built from the word for \textquotedblleft{}country\textquotedblright{}, $\leftarrow$ Late Latin \emph{pagensis}, from \emph{pagus}, \textquotedblleft{}a district\textquotedblright{} (Lewis \& Short), with \emph{-age}.
\end{cousinweb}

\begin{grammarlens}[title={What you've built}]
One more word for this chapter.
\end{grammarlens}

\subsection*{Your turn}
\begin{itemize}
\raggedright
  \item \emph{Say it:} \emph{le paysage}
  \item \emph{Say it:} \emph{le paysage}, once more
  \item \emph{Say it:} say \emph{le paysage}
  \item \emph{Recall:} say the French for departure, then the French for a destination, then say \emph{le paysage} again
\end{itemize}

\begin{tcolorbox}[breakable,colback=teal!4,colframe=teal!35!black,title={Before you move on}]
What does le paysage mean? (\textquotedblleft{}A landscape\textquotedblright{}.) Say it once more.
\end{tcolorbox}
